% Slide theme for the Phase 1 presentation, Suhas H, roll number 25M0328, CSRE, IIT Bombay.
% Replicates the layout of beamer-theme-standard.pdf: white frames, bold blue frame
% titles at the left, the section name at the top, small square bullets; Palatino.
% Three levels, each in its own style: the header (section name, small grey, above a
% hairline), the heading (frame title, large bold blue), the content (black text).
% The footer carries the project title at the left and the two logos at the right.
% Compile from the ppt/ folder of the repository, so that ../figures/ resolves.

% ---------------------------------------------------------------------------
% Type: Palatino, as in the report
% ---------------------------------------------------------------------------
\usepackage[T1]{fontenc}
\usepackage{palatino}
\usefonttheme{serif}
\usefonttheme{professionalfonts}
\usepackage{microtype}

\usepackage{graphicx}
\graphicspath{{../figures/}{../report/}{./}}
\usepackage{booktabs}
\usepackage{array}
\newcolumntype{L}[1]{>{\raggedright\arraybackslash}p{#1}}
\usepackage{tikz}
\usetikzlibrary{positioning,arrows.meta,calc}

% Subscripts in chemical formulae, in the text typeface: SO\sub{2}
\newcommand{\sub}[1]{\textsubscript{#1}}

% ---------------------------------------------------------------------------
% Colours: one blue, the blue of the institute logo, and the theme's orange
% for the rare emphasised word
% ---------------------------------------------------------------------------
\definecolor{iitblue}{HTML}{014EA0}
\definecolor{navy}{HTML}{1A2744}
\definecolor{emph}{HTML}{E08A00}
\definecolor{textcol}{HTML}{111111}
\definecolor{softgrey}{HTML}{6B6B6B}

\setbeamercolor{normal text}{fg=textcol,bg=white}
\setbeamercolor{structure}{fg=iitblue}
\setbeamercolor{frametitle}{fg=iitblue}
\setbeamercolor{title}{fg=iitblue}
\setbeamercolor{item}{fg=iitblue}
\setbeamercolor{alerted text}{fg=emph}
\setbeamercolor{caption name}{fg=iitblue}
\setbeamercolor{block title}{fg=white,bg=iitblue}
\setbeamercolor{block body}{fg=textcol,bg=iitblue!6}

\setbeamerfont{frametitle}{size=\Large,series=\bfseries}
\setbeamerfont{title}{size=\huge,series=\bfseries}
\setbeamerfont{alerted text}{series=\bfseries}
\setbeamerfont{caption}{size=\small}

\setbeamertemplate{navigation symbols}{}
\setbeamertemplate{itemize item}{\small\raise0.25ex\hbox{\vrule width0.55em height0.55em}}
\setbeamertemplate{itemize subitem}{\scriptsize\raise0.3ex\hbox{\vrule width0.5em height0.5em}}
\setbeamertemplate{enumerate item}{\insertenumlabel.}
\setbeamertemplate{blocks}[default]
\setbeamertemplate{caption}[numbered]
\setbeamersize{text margin left=7mm, text margin right=7mm}

% ---------------------------------------------------------------------------
% Header: the current section name, small and grey, at the left above a hairline
% that runs the full text width. It names where the talk is, not the slide.
% ---------------------------------------------------------------------------
\definecolor{hairline}{HTML}{C9CDD4}
\setbeamertemplate{headline}{%
  \vskip2mm
  \hbox to \paperwidth{\hskip7mm
    {\scriptsize\color{softgrey}\insertsectionhead}\hfill}%
  \vskip1.2mm
  \hbox to \paperwidth{\hskip7mm{\color{hairline}\rule{\dimexpr\paperwidth-14mm\relax}{0.4pt}}\hfill}%
}

% ---------------------------------------------------------------------------
% Heading: the frame title, large, bold and blue, at the left. There is no
% subtitle: what the slide shows is said in the content, in black.
% ---------------------------------------------------------------------------
\setbeamertemplate{frametitle}{%
  \vskip3.5mm
  {\usebeamerfont{frametitle}\usebeamercolor[fg]{frametitle}\raggedright
   \insertframetitle\par}%
  \vskip1mm
}

% The one sentence that says what a figure shows: content, so black and upright
\newcommand{\finding}[1]{{\raggedright #1\par}\vskip2mm}

% Literature behind a slide, small and grey at the foot of the content
\newcommand{\refline}[1]{\par\vfill{\raggedright\scriptsize\color{softgrey}#1\par}\vskip0.5mm}

% A small bold blue label that opens a column or a group of points

% Harvey ball: 0 empty, 1 quarter, 2 half, 3 three quarters, 4 full
\newcommand{\hb}[1]{\tikz[baseline=-0.7ex]{%
  \ifcase#1\relax
  \or\fill[iitblue] (0,0) -- (0,1.15mm) arc[start angle=90,end angle=0,radius=1.15mm] -- cycle;%
  \or\fill[iitblue] (0,0) -- (0,1.15mm) arc[start angle=90,end angle=-90,radius=1.15mm] -- cycle;%
  \or\fill[iitblue] (0,0) -- (0,1.15mm) arc[start angle=90,end angle=-180,radius=1.15mm] -- cycle;%
  \or\fill[iitblue] (0,0) circle (1.15mm);%
  \fi
  \draw[iitblue,line width=0.5pt] (0,0) circle (1.15mm);}}
\newcommand{\colhead}[1]{{\bfseries\color{iitblue}#1\par}\vskip1mm}

% Drawn diagrams: flat boxes with a thin blue border, plain arrows
\tikzset{
  box/.style={draw=iitblue,line width=0.6pt,fill=iitblue!5,align=center,
              inner sep=1.8mm,font=\small,text width=#1,minimum height=11mm},
  box/.default={24mm},
  keybox/.style={box=#1,fill=iitblue,text=white,font=\small\bfseries},
  keybox/.default={24mm},
  greybox/.style={box=#1,draw=softgrey,fill=black!4,text=softgrey},
  greybox/.default={24mm},
  flow/.style={-{Stealth[length=2.2mm]},line width=0.7pt,color=iitblue},
  note/.style={font=\scriptsize,text=softgrey,align=center}}

% Outline: numbered sections in blue
\setbeamertemplate{section in toc}{%
  {\color{iitblue}\bfseries\inserttocsectionnumber.}\hskip2mm\inserttocsection\par}
\setbeamerfont{section in toc}{size=\large}

% ---------------------------------------------------------------------------
% Foot: the project title at the left, the two logos small at the right
% ---------------------------------------------------------------------------
\newcommand{\footlogoheight}{4.8mm}
\setbeamertemplate{footline}{%
  \hbox to \paperwidth{%
    \hskip7mm
    \raisebox{1.1mm}{\tiny\bfseries\color{iitblue}\insertshorttitle}%
    \hfill
    \includegraphics[height=\footlogoheight]{iitb-logo}%
    \hskip2.5mm
    \includegraphics[height=\footlogoheight]{csre-logo}%
    \hskip7mm}%
  \vskip2.5mm
}

% ---------------------------------------------------------------------------
% Title page: the two logos at the same height, mirrored about the centre line
% of the page and level with the title; below, the lines of the report's title
% page, one under the other, centred
% ---------------------------------------------------------------------------
% The CSRE mark is about twice as wide as the IIT Bombay seal, so at equal height it
% looks heavier; it is set a little smaller. Each logo is placed from the edge of
% the title itself, at the same gap, so the title sits exactly between them; the
% logos and the title block share one vertical centre.
\newcommand{\titlelogoheight}{13mm}
\newcommand{\csrelogoheight}{10.5mm}
\newcommand{\logogap}{12.8mm}
\newcommand{\talksubtitle}{}
\newcommand{\talkdetails}{}
\newcommand{\maketitleframe}{%
  \begin{frame}[plain,noframenumbering]
    \begin{tikzpicture}[remember picture,overlay]
      \node[anchor=center,align=center,inner sep=0pt] (titleblock)
        at ([yshift=-15mm]current page.north)
        {{\usebeamerfont{title}\usebeamercolor[fg]{title}\inserttitle}\\[1mm]
         {\small\talksubtitle}};
      \node[anchor=east,inner sep=0pt]
        at ([xshift=-\logogap]titleblock.west)
        {\includegraphics[height=\titlelogoheight]{iitb-logo}};
      \node[anchor=west,inner sep=0pt]
        at ([xshift=\logogap]titleblock.east)
        {\includegraphics[height=\csrelogoheight]{csre-logo}};
      \node[anchor=north,text width=0.84\paperwidth,align=center,inner sep=0pt]
        at ([yshift=-33mm]current page.north)
        {\talkdetails};
    \end{tikzpicture}
  \end{frame}
}

% ---------------------------------------------------------------------------
% Closing frame: the words alone, at the exact centre of the page, large, bold, navy
% ---------------------------------------------------------------------------
\newcommand{\makeclosingframe}[1]{%
  \begin{frame}[plain,noframenumbering,c]
    \begin{tikzpicture}[remember picture,overlay]
      \node[anchor=center,inner sep=0pt] at (current page.center)
        {\fontsize{30}{34}\selectfont\bfseries\color{navy}#1};
    \end{tikzpicture}
  \end{frame}
}
